\begin{explanationbox}
\textbf{\textcolor{CExplanation}{一句话直觉}}
互斥事件不会同时发生，所以求它们并集的概率时，可以直接把各自的概率加起来。

\medskip
\textbf{\textcolor{CExplanation}{核心拆解}}
\begin{description}[style=nextline, leftmargin=2.8em, labelsep=0.5em, itemsep=0.45em]
    \item[\textbf{要点一（互斥定义）：}]
    事件$A$与$B$互斥意味着$A\cap B=\varnothing$，即它们没有共同的样本点。

    \item[\textbf{要点二（概率加性）：}]
    对于任意两个事件，都有
    \[
        P(A\cup B)=P(A)+P(B)-P(A\cap B).
    \]
    当$A$与$B$互斥时，$P(A\cap B)=0$，因此
    \[
        P(A\cup B)=P(A)+P(B).
    \]
\end{description}

\medskip
\textbf{\textcolor{CExplanation}{几何本质（维恩图）}}
在维恩图中，互斥事件对应两个没有重叠区域的图形，它们覆盖的总面积就是各自面积之和，类比概率的加法。

\medskip
\textbf{\textcolor{CExplanation}{代数意义（可加性）}}
概率具有可加性：对于互斥事件，其并集的概率等于各事件概率之和；对于有限个两两互斥事件，其并集概率等于各事件概率之和。

\medskip
\textbf{\textcolor{CExplanation}{考点价值（常考方向）}}
\begin{itemize}[leftmargin=*, itemsep=0.5em, topsep=0.4em]
    \item \textbf{考法一（直接计算）：}
    已知$P(A),P(B)$且明确给出互斥，求$P(A\cup B)$。

    \item \textbf{考法二（逆向求值）：}
    在互斥条件下，通过并集概率和其中一个事件概率求另一个。
\end{itemize}

\medskip
\textbf{\textcolor{CExplanation}{顿悟点}}
互斥是“无重叠”，所以加法是“无修正”的加法；一旦出现“至少有一个发生”，先判断是否互斥，再选择公式。

\medskip
\textbf{\textcolor{CExplanation}{使用场景}}
\begin{itemize}[leftmargin=*, itemsep=0.5em, topsep=0.4em]
    \item \textbf{场景一（古典概型）：}
    计算有限等可能样本空间中互斥事件并集的概率。

    \item \textbf{场景二（一般随机试验）：}
    已知事件互斥，计算其并集发生的概率。
\end{itemize}
\end{explanationbox}